\documentclass[11pt,a4paper]{article}
\usepackage[utf8]{inputenc}
\usepackage[vietnamese]{babel}
\usepackage{amsmath,amssymb}
\usepackage{geometry}
\geometry{a4paper, total={180mm,260mm}, left=15mm, right=15mm, top=18mm, bottom=18mm}

\usepackage{tikz}
\usetikzlibrary{shapes.geometric, arrows.meta, positioning, calc}
\tikzset{arr/.style={-{Stealth[scale=0.8]}}}
\usepackage{tcolorbox}
\tcbuselibrary{skins,breakable}
\usepackage{fontawesome}
\usepackage{fancyhdr}
\usepackage{enumitem}

\pagestyle{fancy}
\fancyhf{}
\renewcommand{\headrulewidth}{0.5pt}
\renewcommand{\footrulewidth}{0.5pt}
\fancyhead[L]{\small\bfseries\color{violet!85!black}\faBook\ SỔ TAY KIẾN THỨC TOÁN 2}
\fancyhead[R]{\small\bfseries\color{violet!85!black}Lớp Toán Thầy Lĩnh -- loptoanthaylinh.vn}
\fancyfoot[C]{\small\color{gray}-- Sổ tay kiến thức: Bài 05. Ôn tập phép cộng, phép trừ không nhớ trong phạm vi 100 --}

\setlength{\parindent}{0pt}
\setlength{\parskip}{0pt}

\begin{document}

% --- BANNER BÀI HỌC ---
\begin{tcolorbox}[
  enhanced,
  colback=violet!8!white,
  colframe=violet!70!black,
  arc=3mm, boxrule=1.2pt,
  left=3mm, right=3mm, top=2mm, bottom=2mm,
  halign=center
]
  {\LARGE\bfseries\color{violet!90!black}\faBookmark\ BÀI 05: ÔN TẬP PHÉP CỘNG, PHÉP TRỪ (KHÔNG NHỚ) TRONG PHẠM VI 100}\\[1mm]
  {\itshape\footnotesize\color{black!75} Kĩ thuật đặt tính thẳng cột và tính nhẩm phép cộng, phép trừ không nhớ phạm vi 100 trên 1 trang A4}
\end{tcolorbox}

\vspace{-2mm}

\noindent
\begin{minipage}[t]{0.49\textwidth}\vspace{0pt}
\begin{tcolorbox}[
  enhanced,
  colback=violet!4!white,
  colframe=violet!75!black,
  arc=2mm, boxrule=0.8pt,
  left=2mm, right=2mm, top=1.8mm, bottom=1.8mm,
  title={\footnotesize\bfseries\color{white}\faKey\ 1. Quy tắc đặt tính rồi tính}
]
{\footnotesize
\begin{itemize}\setlength{\itemsep}{1pt}\setlength{\leftskip}{-2.5mm}
  \item \textbf{Bước 1: Đặt tính thẳng cột:}\\
  Viết số hạng này dưới số hạng kia (hoặc số trừ dưới số bị trừ) sao cho:
  \begin{itemize}\setlength{\itemsep}{0.5pt}\setlength{\leftskip}{-2.5mm}
    \item Hàng \textbf{đơn vị} thẳng cột hàng đơn vị.
    \item Hàng \textbf{chục} thẳng cột hàng chục.
  \end{itemize}
  \item \textbf{Bước 2: Thực hiện tính từ phải sang trái:}\\
  Tính từ hàng \textbf{đơn vị trước}, rồi đến hàng \textbf{chục sau}.
  \item Vì là phép tính \textbf{không nhớ}, các chữ số cộng/trừ độc lập từng cột.
\end{itemize}
}
\end{tcolorbox}
\end{minipage}
\hfill
\begin{minipage}[t]{0.49\textwidth}\vspace{0pt}
\begin{tcolorbox}[
  enhanced,
  colback=teal!4!white,
  colframe=teal!75!black,
  arc=2mm, boxrule=0.8pt,
  left=2mm, right=2mm, top=1.8mm, bottom=1.8mm,
  title={\footnotesize\bfseries\color{white}\faSliders\ 2. Kĩ năng tính nhẩm tròn chục}
]
{\footnotesize
\begin{itemize}\setlength{\itemsep}{1pt}\setlength{\leftskip}{-2.5mm}
  \item \textbf{Cộng trừ các số tròn chục:}
  \begin{itemize}\setlength{\itemsep}{0.5pt}\setlength{\leftskip}{-2.5mm}
    \item $40 + 30 = 70$ (Nhẩm: $4\text{ chục} + 3\text{ chục} = 7\text{ chục}$).
    \item $80 - 50 = 30$ (Nhẩm: $8\text{ chục} - 5\text{ chục} = 3\text{ chục}$).
  \end{itemize}
  \item \textbf{Cộng trừ số tròn chục với số có 1 chữ số:}
  \begin{itemize}\setlength{\itemsep}{0.5pt}\setlength{\leftskip}{-2.5mm}
    \item $50 + 6 = 56$ (Gộp $5$ chục và $6$ đơn vị).
    \item $56 - 6 = 50$; $56 - 50 = 6$.
  \end{itemize}
  \item \textbf{Dãy tính 3 số:} Tính lần lượt từ trái sang phải: $40 + 20 - 10 = 60 - 10 = 50$.
\end{itemize}
}
\end{tcolorbox}
\end{minipage}

\vspace{-1.5mm}
\begin{tcolorbox}[
  enhanced,
  colback=orange!4!white,
  colframe=orange!80!black,
  arc=2mm, boxrule=0.8pt,
  left=2mm, right=2mm, top=1.5mm, bottom=1.5mm,
  title={\footnotesize\bfseries\color{white}\faSitemap\ 3. Sơ đồ đối chiếu đặt tính phép cộng và phép trừ không nhớ}
]
\centering
\begin{tikzpicture}[x=1cm, y=1cm]
  % Hộp cộng
  \node[rectangle, rounded corners=2pt, draw=violet!80!black, fill=violet!10, text width=4.5cm, align=left, font=\scriptsize, inner sep=4pt] at (0,0) {
    \textbf{PHÉP CỘNG:} $34 + 25 = 59$\\
    \begin{tabular}{r@{\,}l}
      $3$ & $4$ \\
    $+$ & \\
      $2$ & $5$ \\ \hline
      $\mathbf{5}$ & $\mathbf{9}$
    \end{tabular}
    \quad
    \begin{tabular}{l}
      $\bullet$ $4 + 5 = 9$, viết $9$.\\
      $\bullet$ $3 + 2 = 5$, viết $5$.
    \end{tabular}
  };

  % Hộp trừ
  \node[rectangle, rounded corners=2pt, draw=teal!80!black, fill=teal!10, text width=4.5cm, align=left, font=\scriptsize, inner sep=4pt] at (5.5,0) {
    \textbf{PHÉP TRỪ:} $68 - 24 = 44$\\
    \begin{tabular}{r@{\,}l}
      $6$ & $8$ \\
    $-$ & \\
      $2$ & $4$ \\ \hline
      $\mathbf{4}$ & $\mathbf{4}$
    \end{tabular}
    \quad
    \begin{tabular}{l}
      $\bullet$ $8 - 4 = 4$, viết $4$.\\
      $\bullet$ $6 - 2 = 4$, viết $4$.
    \end{tabular}
  };
\end{tikzpicture}
\end{tcolorbox}

\vspace{-1.5mm}
\noindent
\begin{minipage}[t]{0.49\textwidth}\vspace{0pt}
\begin{tcolorbox}[
  enhanced,
  colback=green!4!white,
  colframe=green!60!black,
  arc=2mm, boxrule=0.8pt,
  left=2mm, right=2mm, top=1.8mm, bottom=1.8mm,
  title={\footnotesize\bfseries\color{white}\faCheckCircle\ 4. Ví dụ thực hành chuẩn mực}
]
{\footnotesize
\textbf{1. Đặt tính rồi tính:}\\
\begin{itemize}\setlength{\itemsep}{0.5pt}\setlength{\leftskip}{-2.5mm}
  \item $42 + 36 = 78$ ($2+6=8$; $4+3=7$).
  \item $79 - 25 = 54$ ($9-5=4$; $7-2=5$).
\end{itemize}
\textbf{2. Tính giá trị biểu thức:}\\
$$85 - 30 + 4 = 55 + 4 = \mathbf{59}$$
$$60 + 20 - 40 = 80 - 40 = \mathbf{40}$$
}
\end{tcolorbox}
\end{minipage}
\hfill
\begin{minipage}[t]{0.49\textwidth}\vspace{0pt}
\begin{tcolorbox}[
  enhanced,
  colback=red!4!white,
  colframe=red!75!black,
  arc=2mm, boxrule=0.8pt,
  left=2mm, right=2mm, top=1.8mm, bottom=1.8mm,
  title={\footnotesize\bfseries\color{white}\faExclamationTriangle\ 5. Mẹo tránh bẫy trắc nghiệm}
]
{\footnotesize
\begin{itemize}\setlength{\itemsep}{1pt}\setlength{\leftskip}{-2.5mm}
  \item \textbf{Bẫy đặt tính lệch hàng:}\\
  Phép tính $54 + 3$: Viết chữ số $3$ thẳng cột với $4$ (hàng đơn vị), \textbf{không} được viết thẳng với $5$.
  \item \textbf{Bẫy tính từ trái sang phải:}\\
  Luôn tạo thói quen tính từ hàng đơn vị trước.
  \item Dãy tính chỉ có phép cộng và trừ thì làm lần lượt từ trái sang phải.
\end{itemize}
}
\end{tcolorbox}
\end{minipage}


\end{document}
